\documentclass[12pt,a4paper]{article}

\usepackage[margin=1in]{geometry}
\usepackage{hyperref}
\usepackage{verbatim}
\usepackage{longtable}
\usepackage{listings}

\lstset{
  basicstyle=\ttfamily\scriptsize,
  breaklines=true,
  breakatwhitespace=false,
  columns=fullflexible,
  keepspaces=true,
  frame=single,
  xleftmargin=0pt,
  xrightmargin=0pt
}

\title{CS-471L Mini Compiler Project Report}
\author{Group: 2023CS32, 2023CS36}
\date{Week 16 Final Deliverable}

\begin{document}
\maketitle

\section{Introduction}

This project implements a mini compiler front end for a compact MicroJava subset. The completed system includes grammar analysis, lexical analysis, top-down parsing, bottom-up LR parsing, semantic analysis through a symbol table, centralized error handling, integration through a single dispatcher executable, and final demonstration outputs.

\section{Grammar Specification}

The source language supports classes, fields, arrays, methods, parameters, local declarations, assignment, calls, conditionals, loops, returns, printing, arithmetic expressions, relational expressions, character constants, and boolean constants.

\subsection{Original and LL(1) Grammar}
\lstinputlisting{docs/grammar.md}

\section{FIRST and FOLLOW Sets}

\lstinputlisting{docs/first_follow.md}

\section{LL(1) Predictive Parsing Table}

The non-recursive predictive parser uses the transformed non-left-recursive grammar and an explicit stack. Its parsing table is included below.

\lstinputlisting{docs/ll1_table.md}

\section{LR Parsing Table Construction}

The bottom-up parser uses the original expression grammar, including left-recursive expression productions. The selected LR method is SLR(1), because the grammar's viable-prefix automaton and FOLLOW-set based reductions are sufficient for the implemented MicroJava subset without requiring canonical LR(1) item expansion.

Construction steps:
\begin{enumerate}
  \item Augment the grammar with a new start production.
  \item Build LR(0) item closures and goto transitions for terminals and non-terminals.
  \item Create shift actions from terminal transitions.
  \item Create reduce actions from completed items using FOLLOW sets of the production left-hand side.
  \item Add accept action for the completed augmented production on end of input.
  \item Store non-terminal transitions in the goto table.
\end{enumerate}

\subsection{LR Action and Goto Tables}
\lstinputlisting{docs/lr_table.md}

\section{Module Descriptions}

\subsection{Lexical Analyzer}
The lexer in \texttt{src/lexer.c} reads the source file with double buffering, tracks line and column numbers, skips whitespace and comments, recognizes keywords, identifiers, numbers, character constants, operators, and punctuation, and returns tokens on demand.

\subsection{Recursive Descent Parser}
The recursive descent parser in \texttt{src/recursive\_descent\_parser.c} implements one routine per grammar non-terminal. It uses the LL(1) transformed grammar and is also the parser connected to semantic analysis.

\subsection{Predictive Parser}
The predictive parser in \texttt{src/predictive\_parser.c} uses an explicit stack and the LL(1) table. It includes synchronization-based recovery for syntax errors.

\subsection{LR Parser}
The LR parser in \texttt{src/lr\_parser.c} uses a state stack, grammar-symbol stack, and the SLR(1) action/goto tables in \texttt{src/lr\_tables.h}. It prints a trace showing stack contents, remaining input, and each shift/reduce/accept/error action.

\subsection{Symbol Table Manager}
The symbol table in \texttt{src/symbol\_table.c} is hash-based and supports nested scopes through a linked scope chain. Each symbol stores name, kind, type, scope level, and declaration line. Insert, lookup, delete, and dump operations are implemented.

\subsection{Error Handler}
The shared error handler in \texttt{src/error\_handler.c} records lexical, syntactic, and semantic errors with line and column numbers, prints clear messages, and provides a final summary.

\section{Integration}

The project can be built as separate lab executables or as one integrated executable. The integrated program is implemented in \texttt{src/compiler.c}. It links the lexer, recursive parser, predictive parser, LR parser, symbol table, and error handler by compiling each parser entry point under a module-specific name and dispatching from a single \texttt{compiler.exe}.

Important commands:
\begin{verbatim}
build.bat
compiler.exe lexer test\samples\hello.mj
compiler.exe recursive test\samples\hello.mj
compiler.exe predictive test\samples\hello.mj
compiler.exe lr test\samples\hello.mj
compiler.exe all test\samples\hello.mj
\end{verbatim}

\section{Sample Inputs}

\subsection{Valid Program: hello.mj}
\lstinputlisting{test/samples/hello.mj}

\subsection{Valid Program: factorial.mj}
\lstinputlisting{test/samples/factorial.mj}

\subsection{Valid Program: arrays\_and\_chars.mj}
\lstinputlisting{test/samples/arrays_and_chars.mj}

\subsection{Invalid Syntax Program}
\lstinputlisting{test/samples/invalid_missing_semicolon.mj}

\subsection{Invalid Semantic Program}
\lstinputlisting{test/samples/invalid_semantic.mj}

\section{Sample Outputs}

\subsection{All Parsers on Valid Input}
\lstinputlisting{output/valid_hello_all.txt}

\subsection{All Parsers on Invalid Syntax Input}
\lstinputlisting{output/invalid_missing_semicolon_all.txt}

\subsection{Semantic Error Output}
\lstinputlisting{output/invalid_semantic_recursive.txt}

\subsection{Symbol Table Dump}
\lstinputlisting{output/symbol_table_hello.txt}

\section{LR Trace}

\lstinputlisting{docs/hello_lr_trace_utf8.txt}

\section{Limitations}

\begin{itemize}
  \item The compiler implements front-end analysis only; it does not generate intermediate code, assembly, or bytecode.
  \item Semantic analysis checks declarations, redeclarations, scope nesting, and identifier usage, but it does not perform full type checking for every expression and function call.
  \item The LR parser reports and recovers from syntax errors by consuming invalid lookahead tokens; recovery is suitable for demonstration but not a full production-grade diagnostic system.
  \item The grammar is a compact MicroJava subset rather than the complete Java language.
\end{itemize}

\section{References}

\begin{itemize}
  \item Aho, Lam, Sethi, and Ullman, \textit{Compilers: Principles, Techniques, and Tools}.
  \item CS-471L Compiler Construction Lab handouts.
  \item Project grammar, parsing tables, and implementation files included in this repository.
\end{itemize}

\section{Conclusion}

All required project sections are complete. The final repository includes grammar artifacts, lexer, three parsers, semantic analysis, error handling, integrated executable support, sample tests, demonstration outputs, and this final report.

\end{document}
